\section{Removing the field and comparing spatial scales}
\label{sec:zero-field}

Proposition~\ref{prop:weak-field} controls an ordinary transverse square in
a small positive field.  We first turn this estimate into an exponential
bound at zero field.  The bound is weak for one prescribed test, but it
holds uniformly over enough tests to compare magnetization averages in a
cell with those in its children.  Throughout this section, $\beta$ is
sufficiently large and fixed.

\subsection{Curvature of a logarithmic Laplace transform}

The analytic step uses the following strip lemma.  Its positivity is
essential: small curvature at the origin will then control curvature on
a real interval whose length grows logarithmically with the spatial
scale.

\begin{lemma}[Strip curvature]\label{lem:strip-curvature}
Let $F$ be holomorphic on $\{z\in\C:|\Ima z|<a_0\}$, real and even on
the real axis, with $F(0)=0$.  Suppose that, for some $K<\infty$,
\[
  0\le F(x)\le K|x|,
  \qquad \Rea F(x+iy)\le F(x)
  \quad(x\in\R,\ |y|<a_0).
\]
For each $a\in(0,a_0)$ there is a strictly positive Schwartz function
$J_a$ such that, with $g_a(x)=F(x)-\Rea F(x+ia)$,
\begin{equation}\label{eq:strip-convolution}
 F''=J_a*g_a,
 \qquad
 J_a=\frac{2}{a^2}(h_a*h_a),
 \qquad
 h_a(x)=\frac{\pi}{2a}\sech^2\!\left(\frac{\pi x}{a}\right).
\end{equation}
In particular, with $c_a=2\pi/a$,
\begin{equation}\label{eq:strip-curvature-growth}
  0\le F''(x)\le e^{c_a|x|}F''(0),
  \qquad
  F(x)\le \frac{F''(0)}{c_a^2}
       \bigl(e^{c_a|x|}-1-c_a|x|\bigr).
\end{equation}
\end{lemma}

\begin{proof}
Put $u(x,y)=\Rea F(x+iy)$.  Schwarz reflection makes $u$ even in $y$.
We first justify Fourier transformation in the unbounded horizontal
direction.  Fix $a'<a_0$.  Around each vertical segment
$\{x+iy:|y|\le a'\}$ choose a rectangle of fixed width and height
strictly inside the strip.  On that rectangle the harmonic function
$M_x-u$, with $M_x=C(1+|x|)$ and $C$ sufficiently large, is positive.
Its value at $(x,0)$ is at most $C(1+|x|)$ because $F(x)\ge0$.
Harnack's inequality, followed by interior derivative estimates, gives
\[
 |u(x,y)|+|\partial_x^j\partial_y^k u(x,y)|
      \le C_{a',j,k}(1+|x|),\qquad |y|\le a',
\]
where for derivatives one first uses a slightly larger closed substrip.
Thus $u(\cdot,y)$ and all derivatives needed below are tempered
distributions, smoothly depending on $y$.

Use the Fourier convention $\widehat h(\xi)=\int e^{-ix\xi}h(x)\dd x$.
The harmonic equation becomes
$\partial_y^2\widehat u=\xi^2\widehat u$.  This equation is solved
\emph{locally in frequency}: multiply by a smooth function of compact
support in $\xi$ and solve the resulting distribution-valued ordinary
differential equation.  Multiplication by its fundamental matrix is
then legitimate, since $\xi$ remains in a compact set.  The initial
data are $\widehat u(\xi,0)=\widehat F(\xi)$ and
$\partial_y\widehat u(\xi,0)=0$.  Consequently
\[
 \widehat u(\xi,a)=\cosh(a\xi)\widehat F(\xi)
\]
on every compact frequency interval.  We do not multiply an arbitrary
tempered distribution globally by the exponentially growing function
$\cosh(a\xi)$.

The smooth function
\[
 m_a(\xi)=\frac{\xi^2}{\cosh(a\xi)-1},
 \qquad m_a(0)=\frac{2}{a^2},
\]
is Schwartz.  The preceding local identity gives
$m_a\widehat g_a=-\xi^2\widehat F$ on every compact frequency interval,
and hence as tempered distributions.  To identify its inverse Fourier
transform, the substitution $r=e^{2\pi x/a}$ and the beta integral give
\[
 \widehat h_a(\xi)
 =\int_0^\infty\frac{r^{-ia\xi/(2\pi)}}{(1+r)^2}\dd r
 =\frac{a\xi/2}{\sinh(a\xi/2)}.
\]
The value at zero is understood by continuity.  It follows that
$m_a=(2/a^2)\widehat h_a^{\,2}$, proving
\eqref{eq:strip-convolution}.  The convolution is an ordinary convergent
integral because $g_a$ has polynomial growth.  Its distributional
identity therefore holds pointwise.

Both $h_a$ and $J_a$ are strictly positive Schwartz functions.  Moreover,
$|h_a'|\le c_a h_a$, so differentiation under the convolution gives
$|J_a'|\le c_a J_a$.  Hence
$J_a(x-t)\le e^{c_a|x|}J_a(-t)$.  Since $g_a\ge0$, integration proves
the first inequality in \eqref{eq:strip-curvature-growth}.  Integrating
twice, using $F(0)=F'(0)=0$ and evenness, proves the second.
\end{proof}

\subsection{A zero-field exponential bound}

\begin{proposition}\label{prop:zero-field}
Fix $S$ and a sufficiently large finite $\beta$.  There are
$R_0=R_0(S,\beta)$, $C_\beta<\infty$, and $\kappa>0$ with the following
property.  Let $L\ge R\ge R_0$, with $L$ even, and partition $\Lambda_L$
into rectangular cells whose side lengths lie in $[R/2,4R]$.  If
$f:\Lambda_L\to[-1,1]$ has sum zero on every cell, then, for
$q=R^{-5/2}$, $b=-\log(1-q)$, and any spin component $a$,
\begin{equation}\label{eq:zero-exp}
 \log\left\langle
   \exp\left(bx\sum_{z\in\Lambda_L}f_zS_z^a\right)
 \right\rangle
 \le C_\beta bV,
 \qquad x\in\R,\quad |x|\le\kappa\log R.
\end{equation}
Consequently there are $A_\beta,c>0$ such that the spectral measurement
of $D_f=\sum_z f_zS_z^a$ in the zero-field state satisfies
\begin{equation}\label{eq:zero-tail}
 \PP\left\{|D_f|>\frac{A_\beta V}{\log R}\right\}
       \le 2e^{-cbV}.
\end{equation}
All bounds are uniform over the indicated partitions and tests.
\end{proposition}

\begin{proof}
It suffices first to take $a=3$ and put the field in direction $Y$.
Write $D_f=\sum_z f_zZ_z$ and define
\[
 F(z)=\frac{1}{bV}\log\langle e^{bzD_f}\rangle_b.
\]
Since $[H_0,M_L^2]=0$, cyclicity of the trace gives
\[
 \langle e^{bzD_f}\rangle_b
 =\frac{\Tr_{\rm phys}
       \bigl(e^{H_0}e^{bM_L^2/2}e^{bzD_f}e^{bM_L^2/2}\bigr)}
       {\Tr_{\rm phys}e^{H_0+bM_L^2}}.
\]
The middle three factors form the restriction of a tensor product of
single-slot sandwich matrices.  Lemma~\ref{lem:strip-maps}, applied to
$e^{bs^2/2}e^{bzf_zs^3}e^{bs^2/2}$, and
Lemma~\ref{lem:trace-zero-free} show that the expression inside the
logarithm has no zeros in a fixed strip $|\Ima z|<a_0$.
Here $R_0$ is large enough that $b\le b_*$ from
Lemma~\ref{lem:strip-maps}.  Choose the analytic logarithm that is real
on the real axis and vanishes at zero.

This is the logarithm of an \emph{ordinary} Laplace transform: the
spectral measure of $D_f$ in the tilted Gibbs state is a probability
measure, even though $D_f$ need not commute with its density matrix.
Rotation by $\pi$ around the field axis makes this measure even.
Therefore $F$ is even, $F(x)\ge0$, and, since $\|D_f\|\le SV$,
\[
 \Rea F(x+iy)\le F(x)\le S|x|.
\]
In particular,
\[
 F''(0)=\frac bV\Var_b(D_f)
       =\frac bV\langle D_f^2\rangle_b
       \le C_\beta R^{-\eta}
\]
by Proposition~\ref{prop:weak-field}, with $\eta=1/40$, and
$b/q\le2$ for large $R$.
This derivative is an ordinary variance, not the Duhamel covariance
obtained by perturbing the Hamiltonian inside its exponential.

Fix $a\in(0,a_0)$ in Lemma~\ref{lem:strip-curvature}.  Choose
$\kappa>0$ so that $c_a\kappa<\eta/2$.  Equation
\eqref{eq:strip-curvature-growth} then gives, uniformly over our tests,
\[
 0\le F(x)\le C_\beta R^{-\eta/2},
 \qquad |x|\le\kappa\log R.
\]
To remove the field, let $\gamma_0$ and $\gamma_b$ be the normalized
zero-field and tilted density matrices.  Since $H_0$ commutes with
$M_L^2$ and $\|M_L^2\|\le SV$,
\[
 \gamma_0\le e^{2bSV}\gamma_b.
\]
Taking the trace against the positive operator $e^{bxD_f}$ costs at
most $2bSV$ in its logarithm and proves \eqref{eq:zero-exp} for $a=3$.
Rotation invariance of $\gamma_0$ gives every common axis.
Finally, exponential Markov at $x=\pm\kappa\log R$, with
$\kappa A_\beta>C_\beta+c$, proves \eqref{eq:zero-tail}.
\end{proof}

\subsection{Comparison through nested partitions}

For a rectangular cell $C$ and an axis $a$, write
\[
 \overline S_C^a=\frac1{|C|}\sum_{x\in C}S_x^a.
\]
All averages in a single fixed axis commute.  The next proposition
compares their ordinary squares; it makes no joint-measurement assertion
about the three spin components.

\begin{proposition}\label{prop:scale-comparison}
Fix $S$ and a sufficiently large finite $\beta$.  For every sufficiently
large fixed $R_*$ and all sufficiently large even $L$, there is a
deterministic rectangular partition $\mathcal P_*(L)$ of $\Lambda_L$
whose side lengths lie in $[R_*/2,4R_*^{10/9}]$ and such that, for
$a=1,2,3$,
\begin{equation}\label{eq:spatial-comparison}
 \frac1V\sum_{C\in\mathcal P_*(L)}|C|
       \left\langle(\overline S_C^a)^2\right\rangle
 \le
       \left\langle(M_L^a/V)^2\right\rangle
       +\frac{C_\beta}{\log R_*}+o_L(1).
\end{equation}
Here $o_L(1)\to0$ at fixed $R_*$ and $\beta$.
\end{proposition}

\begin{proof}
Set $R_0=L$ and $R_{j+1}=R_j^{9/10}$, and stop at the largest $n$
for which $R_n\ge R_*$.  Starting with the one-cell partition, construct
nested rectangular partitions $\mathcal P_j$ with cell sides in
$[R_j/2,4R_j]$.  For example, subdivide each coordinate interval of a
parent into nearly equal integer intervals of length comparable to
$R_{j+1}$.  The ratios $R_j/R_{j+1}$ are large once $R_*$ is large,
so every such subdivision meets the stated bounds.  The final scale
satisfies $R_*\le R_n<R_*^{10/9}$.

Fix an axis and perform its simultaneous site-spin measurement.
Let $D(C)$ denote the measured value of $\overline S_C^a$.  For one
refinement, put $R=R_j$, $r=R_{j+1}$, and define
\[
 T_j=\frac1V\sum_{B\in\mathcal P_j}
             \sum_{\substack{C\in\mathcal P_{j+1}\\C\subset B}}
               |C|\,|D(C)-D(B)|.
\]
To express $T_j$ through the tests of Proposition~\ref{prop:zero-field},
assign a sign $\sigma_C\in\{-1,1\}$ to every child and put
\[
 \bar\sigma_B=\frac1{|B|}\sum_{C\subset B}|C|\sigma_C,
 \qquad
 f_x=\frac{\sigma_C-\bar\sigma_B}{2}
       \quad(x\in C\subset B).
\]
Then $|f_x|\le1$ and $\sum_{x\in B}f_x=0$.  Direct expansion gives
\[
 \sum_x f_xS_x^a
   =\frac12\sum_{B,C\subset B}|C|\sigma_C\bigl(D(C)-D(B)\bigr)
\]
in this measurement.  Maximizing over the signs therefore gives $VT_j/2$.
There are at most $\exp(CV/r^3)$ sign choices, because every child has
volume at least $(r/2)^3$.  Applying \eqref{eq:zero-tail} to each test
and taking a union bound yields
\[
 \PP\{T_j>2A_\beta/\log R\}
 \le 2\exp\{CV R^{-27/10}-cbV\}
 \le 2\exp(-c'VR^{-5/2}).
\]
The last inequality holds uniformly after increasing $R_*$, since
$b\ge R^{-5/2}$ and $R^{-27/10}=o(R^{-5/2})$.

As $|D(C)-D(B)|\le2S$, this also bounds the expected weighted square
increment:
\[
 \E\left[\frac1V\sum_{B,C\subset B}|C|
                 (D(C)-D(B))^2\right]
 \le \frac{C_\beta}{\log R_j}
       +C\exp(-c'VR_j^{-5/2}).
\]
The squares telescope exactly, because $D(B)$ is the weighted mean of
its children's averages:
\[
 \sum_{C\in\mathcal P_{j+1}}|C|D(C)^2
 -\sum_{B\in\mathcal P_j}|B|D(B)^2
 =\sum_{B,C\subset B}|C|(D(C)-D(B))^2.
\]
Now
$\sum_{j<n}(\log R_j)^{-1}\le C/\log R_*$, whereas
$VR_j^{-5/2}\ge L^{1/2}$ and $n=O(\log\log L)$.
Summing the expected increments proves \eqref{eq:spatial-comparison}
with $\mathcal P_*(L)=\mathcal P_n$.  The deterministic construction
works for all three axes.
\end{proof}
